\chapter{Radiation Therapy}

\section*{Overview}
Radiation therapy uses high-energy ionizing radiation (X-rays, gamma rays, or protons) to damage the DNA of cancer cells, preventing them from dividing and growing.

\section*{Alternative Names}
Radiotherapy, Radiation oncological treatment

\section*{Principle}
Ionizing radiation is delivered to a tumor target, damaging the DNA of cancer cells directly and through free radicals. Because dividing cancer cells are more vulnerable to unrepaired DNA damage, they die while surrounding normal tissue, spared by fractionation and precise targeting, recovers.
\section*{History}
Emil Grubbe treated a breast cancer patient with X-rays in 1896, shortly after Roentgen's discovery. Marie Curie pioneered radium treatments in the early 20th century.

\section*{Why It Is Done}
Ordered to cure localized cancers, shrink tumors before surgery, destroy micro-metastases after surgery, or palliate symptoms (pain, bleeding) in advanced cancer.

\section*{Is This a Primary or Secondary Test or Procedure}
Primary curative treatment for many localized cancers (e.g., prostate, head and neck, cervical cancer).

\section*{How It Is Performed}
External beam radiation uses a linear accelerator to direct radiation at the tumor from outside the body. Brachytherapy places radioactive sources directly inside or next to the tumor.

\section*{Preparation}
A simulation session is performed to design custom molds, masks, and marks on your skin to ensure precise alignment for every treatment session.

\section*{Contraindications and Precautions}
Prior maximum radiation tolerance to the same tissue, or active connective tissue diseases (relative, due to severe fibrosis risk).

\section*{Risks}
Skin changes (redness, peeling), fatigue, hair loss in the treated area, and organ-specific side effects (e.g., radiation pneumonitis, cystitis).

\section*{Aftercare and Recovery}
Keep the treated skin clean and dry. Avoid sun exposure. Side effects typically peak 1-2 weeks after treatment ends and then improve.

\section*{Results and Interpretation}
Evaluated over weeks to months using follow-up imaging (CT, PET) to assess tumor response and regression.

\section*{Expected Results}
Reduction or resolution of tumor mass; minimal long-term toxicity to surrounding normal tissues.

\section*{Follow-up}
Regular follow-up with a radiation oncologist, monitoring of side effects, and surveillance scans.

\section*{References}
\begin{enumerate}
  \item MedlinePlus. Radiation Therapy. U.S. National Library of Medicine; 2025.
  \item Current Medical Diagnosis and Treatment. McGraw-Hill; 2025.
\end{enumerate}

\clearpage
